% Revision R1 manuscript — post-dispatch screening/diagnostic framing.
% Journal-agnostic article class; port into the target journal template
% (e.g. IEEE Access) before submission.  Tables live in tables_r1/.
\documentclass[11pt]{article}

\usepackage[margin=1in]{geometry}
\usepackage{amsmath,amssymb,amsthm}
\usepackage{booktabs}
\usepackage{tabularx}
\usepackage{enumitem}
\usepackage{graphicx}
\usepackage[hidelinks]{hyperref}
\usepackage[numbers,sort&compress]{natbib}
\usepackage{siunitx}

\newtheorem{assumption}{Assumption}
\newtheorem{definition}{Definition}
\newtheorem{proposition}{Proposition}
\newtheorem{theorem}{Theorem}
\newtheorem{remark}{Remark}
\newcolumntype{Y}{>{\raggedright\arraybackslash}X}

\newcommand{\R}{\mathbb{R}}
\newcommand{\ones}{\mathbf{1}}
\newcommand{\norm}[1]{\left\lVert #1 \right\rVert}
\newcommand{\abs}[1]{\left\lvert #1 \right\rvert}
\newcommand{\Proj}{\mathrm{Proj}}

\graphicspath{{../figures/}{figures/}}

\title{Adjoint-Based Thermal-Security Robustness Radii for AC Power Systems:\\
A Post-Dispatch Screening and Diagnostic Certificate Under
Response-Constrained Injection Uncertainty}

\author{Egor Grishin\\
Center for Energy Science and Technology, Moscow, Russia\\
\texttt{mkgrgor@gmail.com}}

\date{Revision R1}

\begin{document}
\maketitle

\begin{abstract}
This paper develops a constrained dual-norm framework for first-order
robustness radii of power-system operating points and specializes it to AC
thermal security.  The radius is the largest metric ball of admissible
injection perturbations, restricted by an uncertainty-response model, that
remains feasible for all monitored affine constraint surrogates.  For AC
networks, two-ended apparent-power constraints are mapped to injection space
through sparse transposed-Jacobian solves, avoiding dense sensitivity
matrices; zero-flow line ends, where the apparent-power magnitude is
nondifferentiable, receive an exact first-order operator-norm radius, so the
certificate covers every monitored constraint.  The framework covers
Euclidean, covariance-weighted, and sigma radii, yields worst-case
perturbation directions, and extends to steady-state security-distance
concepts with explicit response constraints.  We prove exact affine-model
maximality, DC slack-reference invariance, and metric--sigma equivalence.
The contribution is positioned strictly as a \emph{post-dispatch screening
and diagnostic layer}: it evaluates a fixed operating point produced by any
dispatch process and certifies the affine surrogate only.  Verification is
correspondingly thorough: the assembled adjoint derivative now agrees with
centered finite differences to the power-flow-tolerance floor (relative
errors decreasing as $O(\varepsilon^2)$ from $10^{-6}$ to $10^{-10}$);
line-wise analytical flow standard deviations match nonlinear Monte Carlo
within $\pm 2\%$ (median ratio 1.004), making the sigma radius a calibrated
first-order index in the tested regime; multi-scale nonlinear replay of
worst-case directions quantifies the mild nonconservatism of the affine
radius (empirical crossing at 0.90--1.00 of the certified radius); and
danger-oriented radius rankings outperform loading- and headroom-only
indicators with scenario-bootstrap confidence intervals that exclude zero.
\end{abstract}

\noindent\textbf{Keywords:} power-system security; security distance;
robustness radius; post-dispatch screening; AC power flow; adjoint
sensitivity; thermal limits; uncertainty quantification.

\section{Introduction}
\label{sec:intro}

Power-system operators need fast, interpretable security margins around an
operating point that has already been selected --- by an OPF, a
security-constrained OPF, a market process, or manual procedures.  The
question this paper answers is strictly \emph{post-dispatch}: given a fixed
AC operating point, how large is the admissible set of continuous injection
perturbations before some monitored thermal constraint is reached, which
constraint is it, and which injection pattern gets there first?

Finite-contingency SCOPF does not answer this continuous, high-dimensional
question, and fast operational metrics --- loading ratio, thermal headroom,
PTDF sensitivities, LODF screening~\citep{WoodWollenbergSheble2013} ---
measure either margin or sensitivity but not a constrained multidimensional
distance to violation.  Robust and chance-constrained OPF internalize
uncertainty during dispatch~\citep{BertsimasSim2004,BenTalNemirovski1998,%
BienstockChertkovHarnett2014,RoaldEtAl2016,LubinBackhausDvorkin2016}, but
they require a chosen uncertainty model and a new dispatch solve, and they do
not evaluate an arbitrary existing operating point.

The idea of measuring security as a \emph{distance} from an operating point
to a security-region boundary is established: steady-state security distance
(SSSD) combines operating margin with boundary sensitivity and has been used
for N$-$1 assessment~\citep{ChenChenXiaKang2013,ChenChenXia2015Concept,%
ChenEtAl2015N1}.  This paper is a constrained-geometry, AC-adjoint member of
that family, not the first distance-to-boundary construction; the precise
positioning is given in Section~\ref{sec:literature} and
Table~\ref{tab:secdist}.  Relative to SSSD, the radius developed here
(i) restricts perturbations to an explicit uncertainty-\emph{response}
subspace $\{\Delta u: A\Delta u = 0\}$ or response map $\Delta u = T\omega$,
so balancing, participation, and reactive-support conventions are part of
the geometry; (ii) supports general metrics $M \succ 0$, recovering
covariance-weighted (sigma) radii as inverse-covariance metric radii;
(iii) evaluates \emph{two-ended AC apparent-power} constraints through
sparse transposed-Jacobian (adjoint) solves, without dense sensitivity
matrices; and (iv) returns closed-form worst-case directions, including an
operator-norm treatment of zero-flow ends where the scalar gradient does not
exist.

The contribution is a unified certificate model and a verified algorithmic
framework:
\begin{enumerate}[leftmargin=*]
  \item a constrained dual-norm radius formula for affine engineering
        constraints, with exact affine-model maximality, worst-case
        directions, DC slack invariance, and metric--sigma equivalence;
  \item DC and AC thermal-security specializations with two-ended
        apparent-power constraints, a reactive-limit-aware active set, and a
        first-order operator-norm certificate for zero-flow ends, so that
        \emph{every} monitored constraint receives a radius;
  \item a sparse AC adjoint implementation whose derivative is verified to
        the power-flow-tolerance floor by step-sweep centered finite
        differences and per-solve adjoint residuals;
  \item a screening-oriented validation methodology: line-wise sigma
        calibration against nonlinear Monte Carlo, multi-scale worst-case
        replay quantifying affine nonconservatism, scenario-bootstrap
        ranking statistics with paired confidence intervals and
        precision/recall, and a physically realizable participation-factor
        response study.
\end{enumerate}

Scope statement, stated once and used throughout: the certificate is
\emph{exact for the affine surrogate} and is a \emph{local first-order
screening indicator} for the nonlinear AC system.  It is not a nonlinear
feasibility certificate, and this paper does not propose a dispatch
optimization method; radius-aware dispatch is left to future work precisely
because a fair evaluation requires matched-cost baselines and a full SCOPF
comparator.

\section{Literature Review}
\label{sec:literature}

\paragraph{Steady-state security distance}
SSSD defines the distance from an operating point to a steady-state
security-region boundary and uses it for assessment and
ranking~\citep{ChenChenXiaKang2013,ChenChenXia2015Concept}; the N$-$1
extension screens contingencies through the same
distance~\citep{ChenEtAl2015N1}.  The present radius differs in the
uncertainty geometry (equality-constrained ellipsoids and covariance metrics
instead of a fixed Euclidean injection norm), the constraint model
(two-ended AC apparent-power limits handled by adjoint solves instead of
DC or boundary-wise linearizations), the explicit worst-direction recovery,
and the computational profile (one sparse LU plus $2m$ triangular solves).
Table~\ref{tab:secdist} gives the side-by-side comparison requested by the
steady-state security-distance literature.

\begin{table}[t]
\centering
\caption{Positioning relative to the steady-state security-distance and
robustness-diagnostics literature.  SSSD = steady-state security distance
\citep{ChenChenXiaKang2013,ChenChenXia2015Concept,ChenEtAl2015N1};
CIA = convex inner approximations of steady-state security regions
\citep{NguyenDvijothamTuritsyn2019}; UD = unified AC operating-point
robustness diagnostics \citep{AntonIlic2026Diagnostics}.
\label{tab:secdist}}
\scriptsize
\begin{tabularx}{\textwidth}{@{}lYYYY@{}}
\toprule
 & SSSD & CIA & UD & This work \\
\midrule
Uncertainty space & bus injections, fixed directions/boundary distance & full nonlinear feasible set & injections + topology, regime changes & response-constrained injections $\{\Delta u: A\Delta u=0\}$ \\
Metric & Euclidean distance to boundary & set inclusion (no scalar metric) & sensitivity-normalized margins & any $M\succ0$: Euclidean, per-unit, covariance ($\Sigma^{-1}$) \\
Constraints & steady-state security region boundary & nonlinear AC feasibility & multiple diagnostic families & two-ended AC apparent-power thermal limits \\
Network model & DC / PV-curve based boundaries & nonlinear AC with bounded remainders & AC sensitivities around solved PF & AC first-order (reduced Jacobian, lossless series model) \\
Contingency treatment & N$-$1 via boundary recomputation & base topology & topological uncertainty classes & DC LODF screening layer (diagnostic only) \\
Nonlinear guarantee & none (boundary linearization) & yes (inner approximation) & none (local) & none beyond affine surrogate; empirically quantified crossing ratios \\
Worst direction & not explicit & not applicable & partially (per diagnostic) & closed form $K_M(A)h/\sqrt{h^\top K_M(A)h}$; operator-norm at $|S_0|{=}0$ \\
Complexity & boundary construction per scenario & SDP/convex program & per-diagnostic linear algebra & one sparse LU + $2m$ triangular solves (chunked) \\
\bottomrule
\end{tabularx}
\end{table}


\paragraph{Nonlinear inner approximations}
Convex inner approximations construct genuinely nonlinear feasibility
certificates by bounding nonlinear terms~\citep{NguyenDvijothamTuritsyn2019}.
They provide stronger guarantees at higher computational cost; the radius
developed here trades the nonlinear guarantee for closed-form evaluation
and scalable adjoint linear algebra, and Section~\ref{sec:results}
quantifies empirically how far the affine radius sits from the nonlinear
crossing distance.

\paragraph{Post-solution diagnostics}
Recent unified diagnostics quantify AC operating-point robustness under
injection and topological uncertainty with regime
changes~\citep{AntonIlic2026Diagnostics}.  That framework spans a broader
diagnostic family (including topology and regime switching); the present
work is narrower and more formulaic: for thermal constraints in a fixed
affine regime it provides a closed-form constrained ellipsoidal radius, an
exact worst-case direction, and a two-ended sparse adjoint implementation
with a machine-verifiable derivative.  The two approaches are complementary:
the radius could serve as one diagnostic inside such a unified suite.

\paragraph{Robust and chance-constrained OPF}
Robust optimization provides feasibility guarantees over uncertainty
sets~\citep{BenTalNemirovski1998,BertsimasSim2004}; chance-constrained OPF
enforces probabilistic security
levels~\citep{BienstockChertkovHarnett2014,RoaldEtAl2016,%
LubinBackhausDvorkin2016}; robust convex restrictions target nonlinear AC
feasibility~\citep{LeeEtAl2021RobustAC}.  These are dispatch-side tools; the
radius is an evaluation-side tool that measures the robustness such methods
achieve, for whatever operating point they produce.

\paragraph{Reliability indices and sensitivity screening}
The sigma radius is structurally a first-order reliability
index~\citep{HasoferLind1974} specialized to affine thermal surrogates with
response-constrained Gaussian geometry; Section~\ref{sec:results} reports
the calibration evidence that justifies (and bounds) this interpretation.
PTDF/LODF and loading-based screening~\citep{WoodWollenbergSheble2013}
remain the fast baselines against which the radius is evaluated.

\section{Problem Formulation}
\label{sec:formulation}

Table~\ref{tab:notation} collects notation and units.

\begin{table}[t]
\centering
\caption{Notation and units.\label{tab:notation}}
\scriptsize
\begin{tabularx}{\textwidth}{@{}llY@{}}
\toprule
Symbol & Units & Meaning \\
\midrule
$\mathcal{N}, \mathcal{L}$ & -- & bus set ($|\mathcal{N}|=n$), monitored line set ($|\mathcal{L}|=m$) \\
$\omega$ & MW, MVAr & primitive uncertainty (forecast errors) \\
$T$ & -- & response map (participation, balancing, coordinate reduction) \\
$\Delta u$ & MW, MVAr & response-coordinate perturbation, $\Delta u = T\omega$ \\
$A$, $C$ & -- & admissibility (equality) matrices in response / primitive space \\
$M$, $W$ & (MW)$^{-2}$ & metric matrices in response / primitive space \\
$\Sigma$ & (MW)$^2$, (MVAr)$^2$ & injection covariance; $M=\Sigma^{-1}$ gives sigma radii \\
$K_M(A)$ & (MW)$^2$ & inverse metric restricted to $\ker A$ \\
$x$ & rad, p.u. & reduced AC state $[\theta_{\mathcal{N}\setminus s}; V_{\mathcal{PQ}}]$ \\
$J$ & MW/rad, MW/p.u. & reduced AC power-flow Jacobian at the base point \\
$b_{\ell,e}$ & MVA/rad, MVA/p.u. & state gradient of $s_{\ell,e}$ \\
$h_{\ell,e}$ & MVA/MW & injection-space sensitivity, $J^\top h_{\ell,e} = b_{\ell,e}$ \\
$s_{\ell,e},\, c_\ell$ & MVA & line-end apparent power; thermal limit \\
$M_{\ell,e} = c_\ell - s^0_{\ell,e}$ & MVA & thermal margin at end $(\ell,e)$ \\
$r^{DC}_\ell$ & MW & DC radius (balanced Euclidean MW injection norm) \\
$r^{AC,M}_{\ell,e}$ & metric-dep. & AC end radius: MW-scale for Euclidean $M$; dimensionless standard deviations for $M=\Sigma^{-1}$ \\
$\beta_j$ & -- & standardized affine margin (sigma radius) of constraint $j$ \\
$\sigma_{\max}$ & MVA/MW & operator norm of the projected 2D $(P,Q)$ response at a zero-flow end \\
\bottomrule
\end{tabularx}
\end{table}


\paragraph{General affine-radius model}
Consider an operating point with monitored affine surrogate constraints
\begin{equation}
  y_j^0+a_j^\top\Delta u \le c_j,
  \qquad j\in\mathcal{J},
\end{equation}
where $\Delta u$ is the local response coordinate.  In primitive space,
uncertainty $\omega$ may pass through a response map $\Delta u=T\omega$
representing balancing participation, storage response, reactive-control
assumptions, or coordinate reduction.  If
$\Omega_W(1)=\{\omega:C\omega=0,\ \omega^\top W\omega\le1\}$, the support
denominator of sensitivity $h$ is
$\sigma(h)=\sqrt{h^\top T K_W(C)T^\top h}$ with $W\succ0$, $C$ full row
rank, and
\begin{equation}
  K_W(C)=W^{-1}-W^{-1}C^\top(CW^{-1}C^\top)^{-1}CW^{-1}.
\end{equation}
Alternatively, work directly in response space with
$\mathcal{B}=\{\Delta u\in\R^d:A\Delta u=0\}$, a metric $M\succ0$, and the
unit set $\mathcal{U}_M(1)=\{\Delta u:A\Delta u=0,\ \Delta u^\top M\Delta
u\le1\}$.  The two descriptions agree when $M$ is the metric induced by
$(T,W,C)$ on $\mathrm{range}(T)$.  For one-sided constraints the margin is
$m_j=c_j-y_j^0$; a two-sided constraint contributes its two one-sided forms.

\begin{definition}[Linearized robustness radius]
$r\ge0$ is a linearized robustness-radius certificate if every
$\Delta u\in r\,\mathcal{U}_M(1)$ satisfies all monitored affine
constraints.  The global radius $r^\star$ is the largest such $r$ when all
base margins are nonnegative.  If any margin is negative, the operating
point is reported as base-infeasible for that constraint rather than
assigned a certified radius.
\end{definition}

\paragraph{Power-system specialization}
Let $\mathcal{N}$ ($|\mathcal{N}|=n$) be the buses and $\mathcal{L}$
($|\mathcal{L}|=m$) the monitored lines with symmetric thermal limits
$c_\ell>0$.  DC certificates monitor active-power flows $f_\ell$ (MW) under
the balanced subspace
$\mathcal{B}_P=\{\Delta p\in\R^n:\ones^\top\Delta p=0\}$.  AC certificates
monitor from- and to-end apparent-power magnitudes $s_{\ell,e}$ (MVA),
$e\in\{f,t\}$, with the response coordinate
$\Delta u=(\Delta P_{\mathcal{I}_P},\Delta Q_{\mathcal{I}_Q})$ on the
retained reduced mismatch coordinates.  The admissibility matrix $A$
encodes the chosen uncertainty-response convention, not a physical law:
examples include net-zero active balance, blockwise active/reactive
net-zero forecast error, or a participation-factor response map
(Section~\ref{sec:results} evaluates one with finite generator headroom).
The radius is exact for the selected affine surrogate; for nonlinear AC
power flow it is a local first-order indicator, and nonlinear replay is an
independent validation layer, not part of the proof.

\paragraph{DC radius}
With DC sensitivities $\Delta f=H\Delta p$, headroom
$M_\ell^{DC}=c_\ell-\abs{f_\ell^0}$, and the ones-complement projection
$\bar g_\ell=g_\ell-(\ones^\top g_\ell/n)\ones$, the line radius is
$r_\ell^{DC}=M_\ell^{DC}/\norm{\bar g_\ell}_2$ (finite case), with
infeasible/inactive classifications as in the implementation.

\paragraph{AC radius}
Let $x=[\theta_{\mathcal{N}\setminus s};V_{\mathcal{PQ}}]$ be the reduced
state and $J=\partial(P,Q)/\partial x$ the reduced Jacobian at the base
point.  \emph{Active-set consistency:} the power flow is solved with
reactive limits enforced; a bus whose every in-service voltage-controlling
generator is saturated at a Q limit is effectively PQ in the converged
solution and is linearized as PQ (its voltage variable and reactive
equation enter $x$ and $J$).  This choice is validated to the
power-flow-tolerance floor in Section~\ref{sec:results}.  For
$s^0_{\ell,e}>0$,
\begin{equation}
  b_{\ell,e}=\frac{P^0_{\ell,e}\nabla_x P_{\ell,e}
  +Q^0_{\ell,e}\nabla_x Q_{\ell,e}}{s^0_{\ell,e}},
  \qquad
  J^\top h_{\ell,e}=b_{\ell,e},
\end{equation}
and the end radius is $r^{AC,M}_{\ell,e}=M^{AC}_{\ell,e}/
\sigma_{\mathcal{U}_M(1)}(h_{\ell,e})$ with
$M^{AC}_{\ell,e}=c_\ell-s^0_{\ell,e}$; the line takes the binding end and
the system takes the minimum line.

\paragraph{Zero-flow ends}
At $s^0_{\ell,e}=0$ the magnitude has no scalar gradient, but the
first-order model is still exact:
$s_{\ell,e}(\Delta u)=\norm{(\nabla P^\top\Delta u,\ \nabla Q^\top\Delta
u)}_2+O(\norm{\Delta u}^2)$.  With $H_{\ell,e}=[h^P_{\ell,e};
h^Q_{\ell,e}]$ obtained from two adjoint solves and projected onto the
admissible subspace, the exact first-order radius is
\begin{equation}
  r_{\ell,e}
  =
  \frac{c_\ell}{\sigma_{\max}\!\left(H_{\ell,e}\,\Proj_{\ker A}\right)},
  \label{eq:opnorm}
\end{equation}
a $2\times2$ eigenvalue computation per end.  The nondifferentiability
threshold is scale-aware, $\max(10^{-9},10^{-9}c_\ell)$ MVA, and
Section~\ref{sec:results} reports its sensitivity.  This closes the gap of
the original submission, where zero-flow ends were excluded and large cases
were only partially certified.

\section{Method}
\label{sec:method}

\paragraph{Constrained support function and exact radius}
With $K_M(A)=M^{-1}-M^{-1}A^\top(AM^{-1}A^\top)^{-1}AM^{-1}$, the support
function of the constrained unit ball is
$\sigma_{\mathcal{U}_M(1)}(h)=\sqrt{h^\top K_M(A)h}$, the worst-case unit
perturbation is $K_M(A)h/\sqrt{h^\top K_M(A)h}$, and:

\begin{theorem}[Exact affine-model radius]
\label{thm:exact-radius}
For one-sided constraints with nonnegative margins $m_j$,
$r^\star=\min_j m_j/\sqrt{a_j^\top K_M(A)a_j}$ is the largest radius for
which every $\Delta u\in r^\star\,\mathcal{U}_M(1)$ satisfies all affine
constraints.
\end{theorem}

The proof is the standard support-function argument for an ellipsoid
intersected with a subspace, specialized to thermal surrogates; following
the convex-geometry literature we present it as a unifying engineering
specialization rather than a new theoretical result.  Slack-reference
invariance of the balanced DC radius and the metric--sigma equivalence
$\beta_j=m_j/\sqrt{h_j^\top K_{\Sigma^{-1}}(A)h_j}$ (the conditional
standard deviation of the affine response under $A\Delta u=0$) carry over
from the original submission unchanged; for a diagonal covariance and a
single balance equation the projection is variance-weighted mean removal.
We emphasize the distinction, raised during review, between
\emph{conditioning} a Gaussian vector on $A\Delta u=0$ and \emph{applying a
physical response map} $T\omega$: the two coincide only when the conditional
covariance equals $T\Sigma_\omega T^\top$, and every experiment below states
which mechanism it uses.

\paragraph{Sparse adjoint algorithm and verification hooks}
The AC computation factorizes the reduced Jacobian once (sparse LU) and
solves one transposed system per line end, in chunks; zero-flow ends add two
solves.  Two verification quantities are computed \emph{inside} the
certificate: the adjoint residual
$\norm{J^\top h-b}_\infty/\max(1,\norm{b}_\infty)$ per chunk (reported with
every run), and the step-sweep centered finite-difference protocol of
Section~\ref{sec:evaluation} (reported with every case study).  Peak
adjoint right-hand-side storage is $O(n_x q)$ at chunk size $q$.

\paragraph{DC N$-$1 screening layer}
For DC contingency screening, LODF-updated sensitivities give
post-contingency radii as a fast diagnostic layer.  This is screening only:
it is not a substitute for AC contingency simulation, and no dispatch
procedure is built on it in this paper.

\section{Verification and Evaluation Protocol}
\label{sec:evaluation}

\paragraph{Benchmarks}
PGLib-OPF networks~\citep{PGLibOPF2019} parsed via
MATPOWER-format inputs~\citep{ZimmermanMurilloSanchezThomas2011} into
pandapower~\citep{ThurnerEtAl2018}; hourly demand profiles derived from
UnitCommitment.jl instances~\citep{UnitCommitmentJL2024} where used.  The
certificate model is the series-only lossless surrogate; the verification
power flow applies the SAME policy, and a regression test asserts equality
of the certificate admittance model and the verification solver's internal
admittance matrix, so the affine model and the replay model cannot drift
apart silently.

\paragraph{Derivative verification}
For each case, the ten tightest differentiable binding ends are replayed
with centered finite differences along the balanced worst-case direction at
$\varepsilon\in\{1,10^{-1},10^{-2},10^{-3}\}$ MVA.  Acceptance requires the
relative error to decrease with $\varepsilon$ until the power-flow-tolerance
floor; a plateau is treated as a model defect, not a diagnostic curiosity.
The original submission's plateau at the $10^{-2}$ level was traced to two
implementation defects --- an ext-grid-inconsistent slack index that shifted
every active-power sensitivity entry by one bus position on affected cases,
and an unmodeled reactive-limit active set --- both fixed and
regression-tested.

\paragraph{Statistical protocol}
Monte Carlo experiments use fixed seeds, report attempted/converged sample
counts and PF failures, and quantify ranking uncertainty by
\emph{scenario-level} bootstrap (resampling Monte Carlo scenarios, which
preserves the cross-line correlation structure), with paired bootstrap
intervals for score differences and precision/recall at $k$ for top-risk
identification.  Timing rows are mean $\pm$ std over 7 repetitions with
measured peak resident memory.

\paragraph{Reproducibility}
All experiments in this revision regenerate from one commit of the public
repository via committed scripts
(\texttt{experiments/revision\_r1/scripts/}), with JSON artifacts committed
alongside (\texttt{experiments/revision\_r1/results/}).  The revision
experiments ran on a Linux container (48-thread Intel Xeon class, Python
3.11, pandapower 2.14, SciPy 1.13, NumPy 1.26); absolute timings therefore
differ from the original submission's hardware, and only relative structure
should be compared.

\section{Results}
\label{sec:results}

\paragraph{Derivative correctness}
Table~\ref{tab:adjoint-validation} shows the step-sweep verification after
the corrections: median relative errors fall from $\sim10^{-6}$ at
$\varepsilon=1$ MVA to $\sim10^{-10}$ at $\varepsilon=10^{-3}$ MVA ---
$O(\varepsilon^2)$ convergence to the power-flow-tolerance floor --- with
adjoint residuals at machine precision.  The non-convergent
$10^{-2}$-plateau of the original submission is fully explained by the two
implementation defects above; no unexplained discrepancy remains.

\begin{table}[t]
\centering
\caption{Adjoint verification after the implementation corrections: centered
finite differences along the balanced worst-case direction of the ten
tightest differentiable line ends per case.  The relative error now decreases
as $O(\varepsilon^2)$ down to the power-flow-tolerance floor, in contrast to
the non-convergent $10^{-2}$-level plateau of the original submission, which
was traced to a slack-index defect and an unmodeled reactive-limit active
set.  Certificate time is mean of 7 repetitions (adjoint + norms, excluding
AC PF).\label{tab:adjoint-validation}}
\scriptsize
\begin{tabularx}{\textwidth}{@{}Xrrrrrrr@{}}
\toprule
Case & Ends & Med.\ rel.\ ($\varepsilon{=}1$) & Med.\ rel.\ ($\varepsilon{=}10^{-2}$) & Med.\ rel.\ ($\varepsilon{=}10^{-3}$) & Max rel.\ ($\varepsilon{=}10^{-3}$) & Adj.\ resid. & Cert.\ (s) \\
\midrule
\texttt{case14\_ieee} & 10 & \num{2.33e-6} & \num{1.98e-10} & \num{1.19e-10} & \num{3.12e-10} & \num{4.06e-16} & 0.0305 \\
\texttt{case30\_ieee} & 10 & \num{6.17e-6} & \num{6.16e-10} & \num{5.36e-11} & \num{6.38e-10} & \num{4.69e-16} & 0.0538 \\
\texttt{case118\_ieee} & 10 & \num{4.11e-7} & \num{3.08e-11} & \num{1.73e-10} & \num{5.51e-10} & \num{4.36e-16} & 0.5022 \\
\texttt{case200\_activ} & 10 & \num{6.62e-9} & \num{7.46e-11} & \num{4.07e-10} & \num{1.68e-9} & \num{1.69e-15} & 0.4795 \\
\bottomrule
\end{tabularx}
\end{table}


\paragraph{DC--AC divergence}
Table~\ref{tab:dc-ac} regenerates the paired radii with the corrected
implementation.  DC radii are bit-identical to the original table, as
implied by slack invariance; AC radii shift moderately where the defects
were active (case73: $14.24\to11.64$; case118: $6.18\to5.80$).  The
qualitative conclusion is unchanged: AC and DC radii differ materially
(AC/DC from 0.45 to 1.16), so DC-only screening misjudges AC thermal
robustness in both directions.

\begin{table}[t]
\centering
\caption{Regenerated paired DC and AC all-constraint radii (ext-grid-consistent
slack, Q-limit-aware active set, operator-norm zero-flow ends; AC certificate
in the balanced two-block $[P;Q]$ Euclidean norm, AC-FPF base points).
\label{tab:dc-ac}}
\footnotesize
\begin{tabularx}{\textwidth}{@{}Xrrrr@{}}
\toprule
Case & $r_\star^{DC}$ & $r_\star^{AC}$ & AC/DC & ND ends \\
\midrule
\texttt{case5\_pjm} & 112.83 & 110.33 & 0.978 & 0 \\
\texttt{case14\_ieee} & 121.42 & 118.88 & 0.979 & 0 \\
\texttt{case24\_ieee\_rts} & 7.90 & 4.20 & 0.532 & 0 \\
\texttt{case30\_ieee} & 12.21 & 8.32 & 0.681 & 0 \\
\texttt{case57\_ieee} & 27.52 & 23.54 & 0.855 & 0 \\
\texttt{case73\_ieee\_rts} & 25.70 & 11.64 & 0.453 & 0 \\
\texttt{case118\_ieee} & 5.02 & 5.80 & 1.156 & 0 \\
\texttt{case200\_activ} & 41.01 & 38.58 & 0.941 & 0 \\
\bottomrule
\end{tabularx}
\end{table}


\paragraph{Complete certification of zero-flow ends}
Table~\ref{tab:zero-flow} demonstrates the operator-norm extension on
\texttt{case2000\_goc}: all 17 lines with zero-flow ends now carry
first-order radii, the all-constraint system radius is defined, and the
nondifferentiable-end count is insensitive to the detection threshold over
eight orders of magnitude.

\begin{table}[t]
\centering
\caption{Zero-flow line ends on \texttt{case2000\_goc} (AC-FPF base point).
Every nondifferentiable end now receives a first-order operator-norm radius,
so the all-constraint certificate is defined; previously these ends were
excluded and the case was only partially certified.  End counts with
$|S_0|$ below a threshold $t$ (MVA): $\le 1e-12$: 42, $\le 1e-09$: 46, $\le 1e-06$: 46, $\le 0.001$: 46, $\le 0.1$: 46 --- the count is insensitive to the
threshold over eight orders of magnitude, supporting the scale-aware default.
\label{tab:zero-flow}}
\footnotesize
\begin{tabularx}{\textwidth}{@{}Xr@{}}
\toprule
Quantity & Value \\
\midrule
Monitored lines & 2743 \\
Lines with a zero-flow end & 17 \\
Operator-norm certified & 17 \\
All-constraint radius defined & yes \\
Global minimum radius & 14.07 \\
Q-limit-saturated buses absorbed as PQ & 19 \\
Max adjoint residual & \num{2.29e-15} \\
\bottomrule
\end{tabularx}
\end{table}


\paragraph{Sigma calibration}
Table~\ref{tab:sigma-calibration} reports line-wise calibration on the IEEE
118-bus system: the analytical first-order standard deviation matches the
nonlinear empirical one with median ratio 1.004 (range 0.989--1.017) over
the ten highest-variance ends, and tightened-limit exceedance probabilities
fall inside or near the Wilson intervals.  The $2.8\times$ miscalibration
reported in the original submission was an artifact of the slack-index
defect, not a property of the method; with the corrected sensitivities the
sigma radius behaves as a calibrated first-order index \emph{in this tested
regime}.  We retain the scope limitation: system-level risk is not the
minimum line-wise index, and calibration under strong nonlinearity, larger
perturbation scales, or heterogeneous covariance families requires the
broader protocol outlined in Section~\ref{sec:discussion}.

\begin{table}[t]
\centering
\caption{Line-wise sigma calibration on the IEEE 118-bus system
(load-proportional $\sigma_P$, $\sigma_Q=\sigma_P\tan(\arccos 0.9)$,
3000 nonlinear Monte Carlo replays, seed 42, 0 PF failures).
Ten highest-variance binding ends.  The analytical first-order standard
deviation now matches the empirical nonlinear standard deviation
(median ratio 1.004, range [0.989, 1.017]);
the tightened-limit exceedance probability at $s_0+2\hat\sigma$ is compared
with the empirical frequency and its Wilson 95\% interval.
\label{tab:sigma-calibration}}
\scriptsize
\begin{tabularx}{\textwidth}{@{}Xlrrrrrrl@{}}
\toprule
Line & End & $|S_0|$ & Ana.\ sd & Emp.\ sd & Ratio & $P_{\mathrm{pred}}$ & $P_{\mathrm{emp}}$ & Wilson 95\% \\
\midrule
\texttt{line\_96} & to & 402.44 & 21.805 & 21.619 & 1.009 & 0.0237 & 0.0193 & [0.0150, 0.0249] \\
\texttt{line\_90} & from & 244.62 & 15.204 & 15.070 & 1.009 & 0.0237 & 0.0223 & [0.0176, 0.0283] \\
\texttt{line\_89} & to & 363.68 & 11.804 & 11.767 & 1.003 & 0.0231 & 0.0253 & [0.0203, 0.0316] \\
\texttt{line\_88} & to & 169.21 & 10.805 & 10.731 & 1.007 & 0.0235 & 0.0200 & [0.0156, 0.0257] \\
\texttt{line\_117} & from & 113.85 & 9.822 & 9.819 & 1.000 & 0.0228 & 0.0343 & [0.0284, 0.0415] \\
\texttt{line\_49} & from & 134.73 & 9.433 & 9.272 & 1.017 & 0.0247 & 0.0270 & [0.0218, 0.0334] \\
\texttt{line\_118} & from & 89.21 & 7.447 & 7.418 & 1.004 & 0.0232 & 0.0250 & [0.0200, 0.0312] \\
\texttt{line\_128} & to & 92.94 & 7.284 & 7.274 & 1.001 & 0.0229 & 0.0200 & [0.0156, 0.0257] \\
\texttt{line\_166} & to & 68.60 & 6.777 & 6.852 & 0.989 & 0.0216 & 0.0217 & [0.0170, 0.0275] \\
\texttt{line\_110} & from & 263.10 & 6.542 & 6.605 & 0.991 & 0.0217 & 0.0250 & [0.0200, 0.0312] \\
\bottomrule
\end{tabularx}
\end{table}


\paragraph{Nonlinear validity range of the affine radius}
Table~\ref{tab:replay} quantifies, per case, where nonlinear violation
actually occurs along the certified worst direction: the interpolated
crossing sits at 0.90--1.00 of the affine radius (median 0.95--0.999).  The
affine radius is therefore mildly \emph{optimistic} --- by at most 10\% and
typically under 5\% on these cases --- and a user who requires conservatism
can de-rate the radius by the reported case-wise factor.  All replays
converged; no voltage-band violations occurred at or below the certified
radius.

\begin{table}[t]
\centering
\caption{Multi-scale nonlinear replay of the balanced worst-case direction
for the five tightest lines per case, at scales
$\alpha\in\{0.25,0.5,0.75,0.9,1.0,1.1,1.25,1.5\}$ of the affine radius.
``Crossing $\alpha$'' is the interpolated scale at which the TARGET line
first violates its limit in the nonlinear replay; values below 1 quantify
the mild nonconservatism of the affine radius along its own worst
direction.  All replays converged; no voltage-band violations occurred at
$\alpha \le 1$.\label{tab:replay}}
\footnotesize
\begin{tabularx}{\textwidth}{@{}Xrrrrrr@{}}
\toprule
Case & Lines & Med.\ crossing $\alpha$ & Min & Max & Violated at $0.9r$ & Q-lim.\ events (base) \\
\midrule
\texttt{case14\_ieee} & 5 & 0.952 & 0.902 & 1.002 & 0/5 & 3 \\
\texttt{case30\_ieee} & 5 & 0.992 & 0.938 & 0.995 & 0/5 & 4 \\
\texttt{case118\_ieee} & 5 & 0.993 & 0.980 & 0.996 & 0/5 & 29 \\
\texttt{case200\_activ} & 5 & 0.999 & 0.958 & 1.000 & 0/5 & 21 \\
\bottomrule
\end{tabularx}
\end{table}


\paragraph{Ranking quality with scenario-level uncertainty}
Table~\ref{tab:ranking} repeats the 118-bus ranking study with 8{,}000
attempted nonlinear scenarios and scenario-bootstrap intervals.  The
radius danger score reaches $\rho=0.781$ against $0.621$ (loading) and
$0.599$ (headroom); the paired differences $+0.159$ and $+0.182$ have
bootstrap intervals bounded away from zero.  The
radius-based danger score dominates loading- and headroom-only scores, and
the \emph{paired} confidence intervals for the Spearman differences exclude
zero, which the line bootstrap of the original submission could not
legitimately claim.  Precision at $k$ and top-$k$ mean overload frequencies
quantify the operational value directly.

\begin{table}[t]
\centering
\caption{Ranking statistics on the IEEE 118-bus system with scenario-level
uncertainty quantification (7612/8000 converged Monte Carlo scenarios,
$\sigma_P=\sigma_Q=30$, seed 42).  Confidence intervals are SCENARIO
bootstrap (resampling Monte Carlo scenarios, 2000 replicates), which
treats the correlated line outcomes correctly, unlike the line bootstrap of
the original submission; replicates re-estimate the empirical frequencies, so
the intervals are mildly attenuated relative to the full-sample point
estimate.  The paired scenario-bootstrap 95\% interval for the
Spearman difference is $\Delta\rho = 0.159$, CI $[0.109, 0.161]$ vs.\ loading and $\Delta\rho = 0.182$, CI $[0.135, 0.182]$ vs.\ headroom.
\label{tab:ranking}}
\footnotesize
\begin{tabularx}{\textwidth}{@{}Xrrrrr@{}}
\toprule
Score & $\rho$ & 95\% CI (scenario) & Prec.@5 & Prec.@10 & Top-5 mean freq. \\
\midrule
AC L2 danger score ($1/r$) & 0.781 & [0.738, 0.774] & 0.60 & 1.00 & 0.878 \\
Loading ratio & 0.621 & [0.599, 0.640] & 1.00 & 1.00 & 0.991 \\
Headroom danger score & 0.599 & [0.574, 0.619] & 0.60 & 1.00 & 0.870 \\
\bottomrule
\end{tabularx}
\end{table}


\paragraph{Realizable response models}
With uncertainty restricted to load buses and generators responding through
headroom-proportional participation factors, the certified radii change by
$-1\%$ to $+21\%$ relative to the blockwise-balanced convention on the five
tightest 118-bus lines, every worst-case perturbation respects all generator
limits, and nonlinear replay at the certified radius reaches
$\abs{S}/c_\ell=1.001$--$1.015$ on the target line --- the certificate is
simultaneously realizable and tight under this response model.  This
supports the central modeling claim: the response convention is part of the
certificate geometry, and changing it changes the certified number in an
auditable way.

\paragraph{Computational profile}
Table~\ref{tab:timing} separates power-flow, operator-build, and certificate
stages with repeated timings and measured peak memory.  The certificate
stage (adjoint solves + norms, including residual checks) remains sub-second
through 200 buses on the revision hardware; the lossless \texttt{runpp}
base point does not converge for \texttt{case2000\_goc}, which is why the
zero-flow study uses the AC feasibility power flow --- the accounting is
reported rather than hidden.

\begin{table}[t]
\centering
\caption{Separated, repeated timings (mean $\pm$ std over 7 repetitions) and
measured peak resident memory.  Stages: nonlinear AC PF; operator build
(Jacobian assembly + sparse LU); certificate (chunked adjoint solves +
balanced norms, including the adjoint-residual check).
\label{tab:timing}}
\scriptsize
\begin{tabularx}{\textwidth}{@{}Xrrrrrrr@{}}
\toprule
Case & Buses & Lines & AC PF (s) & Operator (s) & Certificate (s) & Peak RSS (MB) & Stored $h$ (MB) \\
\midrule
\texttt{case14\_ieee} & 14 & 17 & 0.240$\pm$0.029 & 0.018$\pm$0.005 & 0.030$\pm$0.004 & 323 & 0.01 \\
\texttt{case30\_ieee} & 30 & 34 & 0.275$\pm$0.021 & 0.036$\pm$0.007 & 0.054$\pm$0.005 & 323 & 0.03 \\
\texttt{case118\_ieee} & 118 & 175 & 0.355$\pm$0.021 & 0.129$\pm$0.009 & 0.502$\pm$0.045 & 327 & 0.59 \\
\texttt{case200\_activ} & 200 & 179 & 0.489$\pm$0.147 & 0.186$\pm$0.016 & 0.479$\pm$0.029 & 330 & 1.09 \\
\texttt{case2000\_goc} & \multicolumn{7}{l}{lossless runpp base point does not converge (AC-FPF base required, cf.\ Table~\ref{tab:zero-flow})} \\
\bottomrule
\end{tabularx}
\end{table}


\section{Discussion}
\label{sec:discussion}

\paragraph{What the certificate is}
A post-dispatch screening and diagnostic layer: exact for the affine
surrogate, locally valid for the nonlinear system, calibrated in the tested
regime, and honest about where the nonlinear crossing sits relative to the
certified radius (0.90--1.00 on the studied cases).  Under isotropic
uncertainty several classical scores are monotone transformations of the
same standardized margin; the value of the radius family is that it stays
well-defined and auditable when the geometry is \emph{not} isotropic ---
heterogeneous covariance, response constraints, participation maps.

\paragraph{What it is not}
Not a nonlinear feasibility certificate (unlike convex inner
approximations~\citep{NguyenDvijothamTuritsyn2019}), not a system-level
probability of overload, and not a dispatch method.  Radius-aware dispatch
is deliberately out of scope: a defensible evaluation requires matched-cost
baselines (loading-, headroom-, PTDF-, and chance-based tightening under
equal cost or equal cut budgets) and a recognized preventive/corrective
SCOPF comparator across multiple lossy systems; that is a separate study.

\paragraph{Inequality-constrained uncertainty sets}
Box and polyhedral restrictions (generator headroom, ramp bounds, load
nonnegativity, renewable availability) leave the closed-form family: the
support function becomes a convex program (an SOCP for ellipsoid-plus-box).
The participation-factor study shows one practically relevant middle ground
that stays closed-form; the general trade-off between closed form and
realism is inherent and should be chosen per application.

\paragraph{Limitations}
The AC certificate is local to the linearization at the base point,
including its reactive-limit active set; large perturbations can change the
active set along the path even when the base set is modeled correctly (the
replay protocol reports these events).  The lossless series-only surrogate
is a modeling choice enforced consistently across certificate and
verification; extending the Jacobian to the full pi-model (charging,
shunts, transformer magnetizing branches) removes the need for the lossless
policy and is mechanical but not yet implemented.  Calibration evidence
covers the tested uncertainty scale and one network family; base
infeasibility, PF non-convergence on stressed cases, and islanding
contingencies are reported as explicit statuses rather than silently
dropped.

\section{Conclusions}

This paper provides a response-constrained thermal-security robustness
radius as a verified post-dispatch screening and diagnostic certificate:
a constrained dual-norm formulation with exact affine-model maximality,
slack invariance, and metric--sigma equivalence; a sparse two-ended AC
adjoint implementation whose derivative is machine-verified; operator-norm
radii that certify zero-flow ends and make the all-constraint radius
well-defined on large cases; calibrated sigma diagnostics; quantified
affine-to-nonlinear crossing ratios; statistically defensible ranking
evidence; and a realizable participation-response study.  The framework is
positioned within the steady-state security-distance literature as its
constrained-geometry, AC-adjoint extension, and the dispatch-side use of
the radius is left to a dedicated future study.

\section*{Data and Code Availability}

Source code, entry points, tests, configuration files, revision experiment
scripts, and their JSON artifacts are available in the public repository
\url{https://github.com/nemakgregor/power-stability-radius} (branch
\texttt{revision-r1-screening}; the exact commit hash accompanies each
artifact).  An archived release with a DOI will be minted at
resubmission.  Benchmark cases are public PGLib-OPF
inputs~\citep{PGLibOPF2019}.

\bibliographystyle{unsrtnat}
\bibliography{references_r1}

\end{document}
